\clearpage
\appendix

%%%%%%%%%%%%%%%%%%%%%%%%%%%%%%%%%%%%%%%%%%%%%%%%%%%%%%%%%%%%%%%%%%%%%%%%%%%%%%%%

\section{Polynomials}

%%%%%%%%%%%%%%%%%%%%%%%%%%%%%%%%%%%%%%%%%%%%%%%%%%%%%%%%%%%%%%%%%%%%%%%%%%%%%%%%

The following sections summarize the key properties of the
polynomial functions used in formulating the Flux Reconstruction
method.
%
Additional details can be found in \linecite{Huynh_2007}.

%%%%%%%%%%%%%%%%%%%%%%%%%%%%%%%%%%%%%%%%%%%%%%%%%%%%%%%%%%%%%%%%%%%%%%%%%%%%%%%%

\subsection{Legendre Polynomials}
\label{SEC:LegendrePoly}

%%%%%%%%%%%%%%%%%%%%%%%%%%%%%%%%%%%%%%%%%%%%%%%%%%%%%%%%%%%%%%%%%%%%%%%%%%%%%%%%

The \key*{Legendre polynomial}[polynomial!Legendre]
of degree $k$ on the interval $-1\leq r\leq1$ is given
by the following recurrence relation\cite{Szego_1967,Hildebrand_1974}:
%
\begin{align}
    \Leg{-1}\pvp[r] =& 0 &
    \Leg{ 0}\pvp[r] =& 1 &
    \Leg{ k}\pvp[r] =& \frac{2k-1}{k} r \Leg{k-1}\pvp[r]
                    -  \frac{ k-1}{k}   \Leg{k-2}\pvp[r]
    \quad\text{for $k \geq 1$}
\end{align}
%
Its derivative can likewise be computed from:
\begin{align}
    \dLeg{-1}\pvp[r] =  
    \dLeg{ 0}\pvp[r] =& 0 &
    \dLeg{ k}\pvp[r] =& \frac{2k-1}{k}    \Leg{k-1}\pvp[r]
                      + \frac{2k-1}{k} r \dLeg{k-1}\pvp[r]
                      - \frac{ k-1}{k}   \dLeg{k-2}\pvp[r]
    \quad\text{for $k \geq 1$}
\end{align}
%
Using these formulas, the first few Legendre polynomials and their
derivatives can be written explicitly as:
\begin{align}
     \Leg{0}\pvp[r] =& 1 &
    \dLeg{0}\pvp[r] =& 0 \\
     \Leg{1}\pvp[r] =& r &
    \dLeg{1}\pvp[r] =& 1 \\
     \Leg{2}\pvp[r] =& \left(  3 r^2 -  1         \right) / 2 &
    \dLeg{2}\pvp[r] =& 3 r \\
     \Leg{3}\pvp[r] =& \left(  5 r^3 -  3 r       \right) / 2 &
    \dLeg{3}\pvp[r] =& \left( 15 r^2 -  3         \right) / 2 \\
     \Leg{4}\pvp[r] =& \left( 35 r^4 - 30 r^2 + 3 \right) / 8 &
    \dLeg{4}\pvp[r] =& \left( 35 r^3 - 15 r       \right) / 2
\end{align}
 
%%%%%%%%%%%%%%%%%%%%%%%%%%%%%%%%%%%%%%%%%%%%%%%%%%%%%%%%%%%%%%%%%%%%%%%%%%%%%%%%

These polynomials are plotted in \fref{FIG:LegendrePoly}.
%
%Note that each Legendre polynomial of (even,odd) order is an (even,odd)
%function.
Note that each order of the Legendre polynomial alternates between
being an even or odd function.
%
That is to say,
\begin{gather}
    \Leg{k}\pvp[-r] = (-1)^k \Leg{k}\pvp[r]
    \label{EQN:LegendreEvenOdd}
\end{gather}
%
Values and derivatives on the left and right edges of the cell are given
by\cite{Huynh_2007} (for $k \geq 0$):
\begin{align}
    \Leg{k}\pvp[-1] =& (-1)^k
 &  \Leg{k}\pvp[ 1] =&   1
 \\
    \dLeg{k}\pvp[-1] =& (-1)^{k-1} k (k+1) / 2
 &  \dLeg{k}\pvp[ 1] =&            k (k+1) / 2
\end{align}

%%%%%%%%%%%%%%%%%%%%%%%%%%%%%%%%%%%%%%%%%%%%%%%%%%%%%%%%%%%%%%%%%%%%%%%%%%%%%%%%

The zeros of the Legendre polynomial are referred to as the Gauss points.
%
As seen in the \fref{FIG:LegendrePoly}, the number of such zeros is equal
to the order of the polynomial.
%
%These points are symmetrically, but non-uniformly, distributed across the cell.
These points are symmetrically distributed across the cell, but biased
toward the edges of the cell.
%
For every odd polynomial order, one of the zeros is located in the center
of the cell.

%%%%%%%%%%%%%%%%%%%%%%%%%%%%%%%%%%%%%%%%%%%%%%%%%%%%%%%%%%%%%%%%%%%%%%%%%%%%%%%%

\begin{figure}
  \begin{centering}
    \insertimage[width=4.00in]{Intro/Legendre}
    \caption{Legendre polynomials and Gauss points.}
    \label{FIG:LegendrePoly}
  \end{centering}
\end{figure}

%%%%%%%%%%%%%%%%%%%%%%%%%%%%%%%%%%%%%%%%%%%%%%%%%%%%%%%%%%%%%%%%%%%%%%%%%%%%%%%%

One can also show that the inner product of any two Legendre polynomials
of orders ($k,l$) is given by:
\begin{gather}
    (\Leg{k},\Leg{l})
 =  \int_{-1}^{1} \Leg{k}\pvp[r] \Leg{l}\pvp[r] dr
 =  \begin{cases}
        0               &  \text{for $k \neq l$} \\
        \frac{2}{2k+1}  &  \text{for $k = l$}
    \end{cases}
\end{gather}
%
Therefore, Legendre polynomials of different orders are orthogonal.
%
Also, the inner product of a Legendre polynomial with itself provides a
measure of the magnitude of the polynomial, which is useful in DG schemes
but not of particular importance for FR.

%%%%%%%%%%%%%%%%%%%%%%%%%%%%%%%%%%%%%%%%%%%%%%%%%%%%%%%%%%%%%%%%%%%%%%%%%%%%%%%%

Using the above relation, one can further show that for $k \geq 1$:
\begin{gather}
    (\Leg{0},\Leg{k})
 =  \int_{-1}^{1} \Leg{k}\pvp[r] dr
 =  0
\end{gather}
This indicates that the net area under any Legendre polynomial curve
of degree $\geq 1$ is zero.

%%%%%%%%%%%%%%%%%%%%%%%%%%%%%%%%%%%%%%%%%%%%%%%%%%%%%%%%%%%%%%%%%%%%%%%%%%%%%%%%

\subsection{Radau Polynomials}
\label{SEC:RadauPoly}

%%%%%%%%%%%%%%%%%%%%%%%%%%%%%%%%%%%%%%%%%%%%%%%%%%%%%%%%%%%%%%%%%%%%%%%%%%%%%%%%

The left and right \key*{Radau polynomials}[polynomial!Radau]
are defined in terms of \key{Legendre polynomials}[polynomial!Legendre].
\begin{align}
    \RadL{k}\pvp[r] =& \frac{  1   }{2} \left( \Leg{k  }\pvp[r]
                                             + \Leg{k-1}\pvp[r] \right)
 &  \RadL{k}\pvp[-1] =& 0
 &  \RadL{k}\pvp[ 1] =& 1
 \nonumber % It is easier to ref this whole group as 1 eqn, not 2.
 \\
    \RadR{k}\pvp[r] =& \frac{(-1)^k}{2} \left( \Leg{k  }\pvp[r]
                                             - \Leg{k-1}\pvp[r] \right)
    \label{EQN:RadauPoly}
 &  \RadR{k}\pvp[-1] =& 1
 &  \RadR{k}\pvp[ 1] =& 0
\end{align}
and are shown in \fref{FIG:RadauPoly}.
%
From \eref{EQN:LegendreEvenOdd},
\eref{EQN:RadauPoly},
and  \fref{FIG:RadauPoly},
one can see they are mirror images of each other.
\begin{gather}
   %\RadR{k}\pvp[-r] = (-1)^k \RadL{k}\pvp[r]
    \RadR{k}\pvp[-r] = \RadL{k}\pvp[r]
\end{gather}
We therefore focus on the left Radau polynomials,
the first few of which and their derivatives can be written as:
\begin{align}
     \RadL{1}\pvp[r] =& \left( r + 1 \right) / 2
 &  \dRadL{1}\pvp[r] =& 1 / 2
 \\  \RadL{2}\pvp[r] =& \left(                    3 r^2 +  2 r - 1 \right) / 4
 &  \dRadL{2}\pvp[r] =& \left(                    3 r   +  1       \right) / 2
 \\  \RadL{3}\pvp[r] =& \left(           5 r^3 +  3 r^2 -  3 r - 1 \right) / 4
 &  \dRadL{3}\pvp[r] =& \left(          15 r^2 +  6 r   -  3       \right) / 4
 \\  \RadL{4}\pvp[r] =& \left( 35 r^4 + 20 r^3 - 30 r^2 - 12 r + 3 \right) / 16
 &  \dRadL{4}\pvp[r] =& \left( 35 r^3 + 15 r^2 - 15 r   -  6       \right) / 4
\end{align}



